\documentclass[11pt]{article}
\usepackage[a4paper,margin=2.3cm,headheight=22pt]{geometry}
\usepackage[T1]{fontenc}
\usepackage{iftex}
\ifPDFTeX\usepackage[utf8]{inputenc}\fi
\usepackage{lmodern,microtype,booktabs,longtable,tabularx,array}
\usepackage{enumitem,xcolor,colortbl,titlesec}
\usepackage{xurl,hyperref,fancyhdr}
\definecolor{ink}{HTML}{183247}
\definecolor{accent}{HTML}{126C70}
\definecolor{muted}{HTML}{526676}
\definecolor{paperblue}{HTML}{EDF6F5}
\definecolor{gold}{HTML}{B68032}
\definecolor{papergold}{HTML}{FBF4E8}
\definecolor{rulecolor}{HTML}{B8D4D3}
\hypersetup{colorlinks=true,linkcolor=accent,urlcolor=accent}
\titleformat{\section}
  {\Large\bfseries\sffamily\color{ink}}
  {\colorbox{accent}{\color{white}\strut\thesection}}{0.6em}{}
  [\vspace{2pt}{\color{gold}\titlerule[0.7pt]}]
\titleformat{\subsection}
  {\large\bfseries\sffamily\color{accent}}{\thesubsection}{0.6em}{}
\titlespacing*{\section}{0pt}{12pt}{8pt}
\titlespacing*{\subsection}{0pt}{10pt}{5pt}
\arrayrulecolor{rulecolor}
\newcommand{\tablehead}{\rowcolor{paperblue}}
\setlength{\emergencystretch}{4em}
\setlength{\parindent}{0pt}
\setlength{\parskip}{6pt}
\setlist{itemsep=3pt,topsep=4pt}
\setcounter{tocdepth}{2}
\raggedbottom
\newcolumntype{L}[1]{>{\raggedright\arraybackslash}p{#1}}
\pagestyle{fancy}
\fancyhf{}
\lhead{\small\sffamily\textcolor{accent}{\textbf{CAMFRANGLAIS}}\;
  \textcolor{muted}{/ COMPILER CONSTRUCTION}}
\rhead{\small\sffamily\textcolor{muted}{CS4110 / SET A}}
\cfoot{\colorbox{paperblue}{\makebox[1.8em]{\small\sffamily\textcolor{accent}{\thepage}}}}
\renewcommand{\headrulewidth}{0.6pt}
\renewcommand{\headrule}{\color{accent}\hrule width\headwidth height\headrulewidth\vskip-\headrulewidth}
\newcommand{\notice}[2]{%
  \par\medskip\noindent\begingroup\setlength{\fboxsep}{8pt}%
  \colorbox{paperblue}{%
    \begin{minipage}{\dimexpr\linewidth-2\fboxsep\relax}
      {\bfseries\sffamily\textcolor{accent}{#1}}\par
      #2
    \end{minipage}}\endgroup\par\medskip}
\newcommand{\projectversion}{4.1}
\newcommand{\projectdate}{27 September 2026}
\newcommand{\projectcoverlabel}{DETERMINISTIC / PUBLIC ACCESS / EVIDENCE-LED}
\newcommand{\projectcover}[3]{%
  \hypersetup{pageanchor=false}
  \begin{titlepage}
    \noindent\colorbox{ink}{%
      \begin{minipage}{\dimexpr\linewidth-2\fboxsep\relax}
        \vspace{8pt}\hspace{8pt}%
        {\small\bfseries\sffamily\color{white}CS4110 \quad / \quad SUMMER 2026 \quad / \quad SET A}
        \par\vspace{8pt}
      \end{minipage}}\par
    \vspace{1.0cm}
    {\color{gold}\rule{3.4cm}{4pt}}\par
    \vspace{0.6cm}
    {\Huge\bfseries\color{ink}#1\par}
    \vspace{0.9cm}
    {\LARGE\sffamily\color{accent}Camfranglais\par}
    \vspace{0.4cm}
    {\large\color{muted}#2\par}
    \vspace{1cm}
    \notice{\projectcoverlabel}{#3}
    \vfill
    {\small\bfseries\sffamily\color{accent}AUTHOR GROUP}\par
    {\renewcommand{\arraystretch}{1.45}
    \begin{tabularx}{\linewidth}{@{}X r@{}}
      \toprule\tablehead\textbf{Name} & \textbf{Matricule}\\\midrule
      Kwete Ngouba Junior Rayan & ICTU20241377\\
      Djemtchimo Noukui Bruno Jonatan & ICTU20241585\\
      Amina Boubakary & ICTU20241870\\
      \bottomrule
    \end{tabularx}}
    \par\vspace{0.7cm}
    {\small\textcolor{accent}{\textbf{Version \projectversion}} \quad|\quad \projectdate\par
    Faculty of Information and Communication Technologies\par
    Instructor: Engr. Tanwi Nkiamboh\par
    Assignment due: 29 September 2026}
  \end{titlepage}
  \hypersetup{pageanchor=true}}

\usepackage{graphicx}
\hypersetup{
  pdftitle={Camfranglais - Two Automata Explained Simply},
  pdfauthor={Kwete Ngouba Junior Rayan; Djemtchimo Noukui Bruno Jonatan; Amina Boubakary}
}
\setlength{\parskip}{5pt}
\setlist{itemsep=2pt,topsep=3pt}
\newcommand{\code}[1]{\texttt{#1}}
\newcommand{\automaton}[3]{%
  \par\smallskip
  {\centering
    \includegraphics[width=\linewidth,height=#2,keepaspectratio]{diagrams/#1.png}
    \par}
  {\small #3\par}}
\begin{document}

{\LARGE\bfseries\sffamily\color{ink}Our two automata, made simple\par}
An \textbf{automaton} is a machine that follows rules to read input.

\section{The lexer: finding the pieces}
The lexer cuts text into \textbf{tokens}: words, numbers and punctuation.
Its \textbf{DFA} (deterministic finite automaton) tracks the current state;
it has no stack.

\automaton{automaton-lexer}{6.3cm}{
\textbf{Reading the arrows:} C = a letter-like word character;
D = a digit; M = a combining accent mark.
A \textbf{double circle} means a complete token can end here.}

\subsection*{How it works}
\begin{enumerate}
\item \textbf{Start at q0.} A letter leads to W (word), a digit to I (integer),
punctuation to P, and another non-space character to O.
\item \textbf{Follow the arrows.} Stay in W for more letters, or I/F for
more digits. J waits for a letter after an apostrophe or hyphen.
R waits for a digit after a decimal point or comma.
\item \textbf{Keep the longest complete token and restart.} Skip spaces.
J and R are waiting states, so a token cannot finish there.
\end{enumerate}

\subsection*{Three small examples}
\begin{tabularx}{\linewidth}{@{}L{2.0cm}X@{}}
\toprule\tablehead\textbf{Input} & \textbf{What happens}\\\midrule
\code{j'ai} & j: W; apostrophe: J; a and i: W. Output: one word.\\
\code{3.14} & 3: I; dot: R; 1 and 4: F. Output: one number.\\
\code{3.} & No digit follows the dot. Output: number \code{3}, then punctuation \code{.}\\
\bottomrule
\end{tabularx}

\notice{What goes to the parser?}{
After splitting \code{je suis kass}, vocabulary classification labels the
three tokens \textbf{French function word}, \textbf{verb}, \textbf{adjective}.
The parser receives these categories.}

{\small Python regular expressions perform the recognition shown above.
Classification is a separate step: a word-shaped token can still be
\code{UNKNOWN}.\par
\textbf{Main files:} \nolinkurl{compiler/lexer/tokenizer.py} and
\nolinkurl{compiler/lexer/lexicon.py}.}

\clearpage
\section{The parser: checking the order}
The \textbf{pushdown automaton} adds a \textbf{stack}: a last-in, first-out
to-do list. Our \textbf{LL(1) parser} reads categories from left to right,
using one category of lookahead to choose a grammar rule.

\automaton{automaton-parser}{4.2cm}{
\textbf{States:} q0 starts; q\_load gets the next category;
q[a] holds category a while the parser works.
\textbf{\$} marks the end; \textbf{eps} means no input is read on that arrow.}

\subsection*{The three main actions}
\begin{enumerate}
\item \textbf{Expand.} Start with \code{[\$, Ut]} (Ut = whole utterance).
Replace a grammar symbol with its table rule, pushed backwards so the
first symbol is on top. Do not read a token yet.
\item \textbf{Match.} Compare the expected category on top with the current
input category. If equal, remove the top and move to the next token.
\item \textbf{Finish.} Empty rules just remove a grammar symbol.
Accept when only \code{\$} remains and input has ended. A mismatch or
missing rule means reject.
\end{enumerate}

{\small\textbf{Diagram notation:} \code{M[A,a]} selects a rule for grammar
symbol A and category a; \code{rev(rhs)} pushes its replacement backwards.
An arrow reads ``input, stack top / replacement'': \code{a, X / X} reads a
and keeps X. The \textbf{rightmost} stack symbol is the top.}

\subsection*{A simple successful run}
For \code{je suis kass} (``I am tired''), the categories are
\code{FF, V, ADJ}: French function word, verb, adjective.
We group the intermediate grammar expansions below.

\begin{tabularx}{\linewidth}{@{}L{2.8cm}X@{}}
\toprule\tablehead\textbf{Stage} & \textbf{In simple language}\\\midrule
Prepare & Put Ut on the stack and apply rules until FF is expected.\\
Read \code{je} & Match FF. The next rules prepare for a verb.\\
Read \code{suis} & Match V. The next rule allows an adjective.\\
Read \code{kass} & Match ADJ. There are no more input tokens.\\
Finish & Remove empty-rule symbols. End markers agree: \textbf{accept}.\\
\bottomrule
\end{tabularx}

\textbf{Rejected example:} \code{je suis alli} stops at token 3:
\code{alli} is \code{UNKNOWN}, with no applicable grammar rule.
It is not skipped.

\notice{The difference to remember}{
\textbf{Lexer:} ``What are the pieces?''\quad
\textbf{Parser:} ``Do those pieces fit our grammar?''\par
Accept means grammar fit, not a guarantee of meaning.}

{\small\textbf{Main files:} \nolinkurl{compiler/parser/predictive.py}
(stack machine), \nolinkurl{compiler/parser/yaounde.py} (corpus grammar).}

\end{document}
